\scanpage{259}
\mspage{133}

La démonstration est celle donnée \deleted{pour} 9.7.
on prouve que si $\unsure{X|I},\unsure{Y|J}$ sont deux objets
de $\mathrm{Cat}$, \inserted{\uline{non vides}}, alors deux flèches
$f_0,f_1:X\mathrel{\unsure{=|\rightrightarrows}}Y$ sont égales dans $\mathrm{Hot}_{W}$,
par un argument d'homotopie, en
choisissant deux pts $i_0,i_1$ de $I$ qui
sont dans deux comp. connexes distinctes,
et tels que leurs images \unsure{$f_0,f_1$|$fi_0,fi_1$} dans $J$
soient dans la $\unsure{\overline{\mathrm m}|=}$ composante.

\uline{Prop. 9.11}. \deleted{La conjonction des} \inserted{Soit $W$ un localiseur fondamental}
(i.e. on a $\mathrm{Loc}(1,\unsure{2|1},3)$). Alors les conditions
suivantes sont équivalentes

a) On a $\mathrm{Loc}(5\deleted{\unsure{a}}\inserted{\mathrm{bis}})$ (axiome des \deleted{\unsure{diviseurs}}\inserted{pointes},
forme \deleted{\unsure{faible}}\inserted{forte}) et $\mathrm{Loc}(8)$ (axiome de connexité)

b) On a $\mathrm{Loc}.(5\,\mathrm{bis})$, et $\mathrm{Hot}_{W}$ n'est pas
équivalente à $e$ ou à $\Delta^1$ (cf. commentaire
9.8.)

c) Pour tt $f:X\to Y$, on a $f\in W$ sss
$\pi_0(f)$ iso, et \unsure{pour|après} \unsure{toute|tte} comp. \unsure{connexe|connexes}
$\unsure{X_i|X'}$ de $X$, si $Y_i$ désigne la comp correspondante de $Y$, \unsure{la flèche induite|les flèches induites} $f_i:X_i\to Y_i$
est dans $W$.

\deleted{De plus, ces conditions impliquent}

La raison \deleted{principale} pour laquelle je n'érige
pas cette propriété en un des \unsure{-} axiome
$\mathrm{Loc}(\unsure{\ldots|---})$, c'est que ce n'est \uline{pas} un axiome
de stabilité, mais un axiome de nature
\unsure{mixte}.
% archive sheet number: 258
